\documentclass[10pt,a4paper]{amsart}
\usepackage[margin=19mm]{geometry}
\usepackage{amsmath,amssymb,mathtools}
\usepackage[scr=rsfs,cal=euler]{mathalfa}
\usepackage{xfrac}
\usepackage[unicode,hidelinks,bookmarks=false]{hyperref}
\newcommand{\ie}{i.e.\ }
\newcommand{\cf}{cf.\ }
\newcommand{\resp}{respectively\ }
\newcommand{\Subsection}{\subsection}
\providecommand{\nobreakdash}{\nobreak\hbox{-}}
\setlength{\emergencystretch}{2em}
\multlinegap0pt
\usepackage{bm}
\usepackage{bbm}
\usepackage{dsfont}
\usepackage{xspace}
\usepackage{cancel}
\usepackage{mathtools}
\usepackage[shortlabels]{enumitem}
\usepackage{amsthm}
\makeatletter
\@ifundefined{theorem}{\newtheorem{theorem}{Theorem}}{}
\@ifundefined{lemma}{\newtheorem{lemma}[theorem]{Lemma}}{}
\makeatother
\newcommand{\C}{\mathbb{C}}
\newcommand{\Q}{\mathbb{Q}}
\newcommand{\R}{\mathbb{R}}
\title[Module Lattice Security (Part II): Module Lattice Reduction ]{Module Lattice Security (Part II): Module Lattice Reduction via Optimal Sign Selection}
\author{Ming-Xing Luo}
\date{Source: arXiv:2604.22900v2 (CC BY 4.0)}
\begin{document}
\maketitle

\noindent\emph{Source and licence.} Module Lattice Security (Part II): Module Lattice Reduction via Optimal Sign Selection. Version arXiv:2604.22900v2 (CC BY 4.0), distributed under the Creative Commons Attribution 4.0 International licence, as recorded in the arXiv licence metadata for this exact version and shown on the abstract page https://arxiv.org/abs/2604.22900v2. Full rights notice: licenses/arxiv-2604.22900-CC-BY-4.0.txt.\par\smallskip

\noindent\textit{Editorial scope.} Counted unit: the Lemma whose source label is \texttt{lem:cover} in the article (source0.tex lines 373--379, source label \texttt{lem:cover}), delivered together with the article's own complete original proof of it (source0.tex lines 381--419, one \texttt{proof} environment of 39 lines). No mathematics is written, repaired or re-derived: statement and proof are the article's own text line for line. Required context retained uncounted: lem:orth (lines 318--326). Every definition, display and label of those lines is retained unchanged. Omitted from this package and not redistributed: 1--317, 327--372, 380--380, 420--1778. Nothing inside a retained excerpt is omitted or altered: every retained body is delivered exactly as the article's lines are written, so no omission receipt of any kind accompanies this package. Citations: the retained excerpts carry no citation command at all, so no bibliography entry of the article is transcribed and no reference is redistributed. Reference adaptation: none was needed and none was made; every reference inside the delivered unit resolves inside it, so no unresolved reference or citation is produced. Printed slips kept literal: no printed word, symbol, hypothesis or number of the counted statement or of its proof was altered. Layout and class: the article's own class is \texttt{article} with packages including geometry, amssymb, amsmath, mathtools, amsthm, amsfonts, bbm, enumitem, booktabs, siunitx, fontawesome5, algorithm, algpseudocode, tikz; the wrapper is the lane's pinned \texttt{amsart} editorial wrapper at 10pt on A4, which restates the theorem-like environments the retained text needs and adds the article's own macro definitions that the retained text needs (\texttt{\textbackslash C}, \texttt{\textbackslash Q}, \texttt{\textbackslash R}), with extra wrapper packages (bm, bbm, dsfont, xspace, cancel, mathtools, enumitem). The delivered numbering is the unit's own. Assets: no retained span contains a figure environment, a graphics-inclusion command, a PDF-page inclusion, a table or a TikZ picture; no image file is copied into this package, the package declares no asset dependency and no image file is redistributed with it. Licence: version arXiv:2604.22900v2 of this article is distributed under the Creative Commons Attribution 4.0 International licence, as recorded in the arXiv licence metadata for this exact version and shown on the abstract page https://arxiv.org/abs/2604.22900v2. A full rights notice is delivered at licenses/arxiv-2604.22900-CC-BY-4.0.txt. Characters: every retained excerpt is pure ASCII, so the delivered engine, XeLaTeX with its default Latin Modern text font, reports no missing character.
\begin{lemma}\label{lem:orth}
For $0\le a, b < n$, we have the trace orthogonality 
\begin{align}
  \langle\sigma(\zeta^a),\sigma(\zeta^b)\rangle=n \delta_{a,b},
  \label{E:traceorth}
\end{align}
where $\delta_{a,b}$ denotes the delta function.

\end{lemma}

\begin{lemma}\label{lem:cover}
For any $t\in K\otimes_\Q\R$, the following rounding bounds hold:
\begin{enumerate}[label=(\roman*),nosep]
\item There exists $c\in R$ such that $\|\sigma(t-c)\|_\infty \le \frac{n}{2}$.
\item The same $c$ satisfies $\|\sigma(t-c)\|_2 \le \frac{n}{2}$.
\end{enumerate}
\end{lemma}

\begin{proof}
Since $\{1, \zeta, \dots, \zeta^{n-1}\}$ is a $\Q$-basis of $K$, any element $t \in K\otimes_\Q\R$ admits a unique real-coefficient expansion as
\begin{align}
t=\sum_{m=0}^{n-1} x_m \zeta^m, \ x_m \in \R.
\label{E:exp}
\end{align}
For each $m$, let $c_m=\lfloor x_m \rceil$ denote the nearest integer to $x_m$. Define $c=\sum_{m=0}^{n-1} c_m \zeta^m \in R$. The rounding error coefficients $\epsilon_m=x_m-c_m$ satisfy $|\epsilon_m| \le 1/2$ for all $m$.

(i) For any Galois embedding $\sigma_\ell: K \hookrightarrow \C$ (corresponding to odd $\ell$ modulo $2^k$), we have
\begin{align}
\sigma_\ell(t-c)=\sum_{m=0}^{n-1} \epsilon_m \sigma_\ell(\zeta^m).
\label{E:emb}
\end{align}
Since $\sigma_\ell(\zeta^m)=\zeta^{\ell m}$ is a root of unity, it implies $|\sigma_\ell(\zeta^m)|=1$ for all $m$ and $\ell$. Using the triangle inequality $|x+y|\leq |x|+|y|$, we get 
\begin{align}
|\sigma_\ell(t-c)| \le \sum_{m=0}^{n-1} |\epsilon_m| \cdot |\sigma_\ell(\zeta^m)|=\sum_{m=0}^{n-1} |\epsilon_m|\le \frac{n}{2}
\label{E:sigtri}
\end{align}
as each term is bounded by $1/2$. Since this holds for all embeddings $\sigma_\ell$'s, we get $\|\sigma(t-c)\|_\infty \le n/2$.

(ii) From Lemma \ref{lem:orth}, the embedding vectors $\{\sigma(\zeta^m)\}_{m=0}^{n-1}$ are pairwise orthogonal with squared norm $n$, i.e., 
\begin{align}
\langle \sigma(\zeta^a), \sigma(\zeta^b) \rangle=n\cdot \delta_{a,b}.
\label{E:typair}
\end{align}
We then obtain the squared Euclidean norm as
\begin{align}
\|\sigma(t-c)\|_2^2&=\Big\langle \sum_{m=0}^{n-1} \epsilon_m \sigma(\zeta^m), \sum_{m=0}^{n-1} \epsilon_m \sigma(\zeta^m) \Big\rangle
\notag
  \\
&=\sum_{a,b=0}^{n-1} \epsilon_a \epsilon_b \langle \sigma(\zeta^a), \sigma(\zeta^b) \rangle=n \sum_{m=0}^{n-1} \epsilon_m^2.
\end{align}
Since $|\epsilon_m| \le 1/2$, we get $\epsilon_m^2 \le 1/4$ for all $m$. This further implies that
\begin{align}
\|\sigma(t-c)\|_2^2 \le \frac{n^2}{4}.
\label{E:signorm2}
\end{align}
It yields the bound $\|\sigma(t-c)\|_2 \le n/2$.
\end{proof}

\end{document}
